\documentclass[11pt,a4paper]{article}
\usepackage[T1]{fontenc}
\usepackage[utf8]{inputenc}
\usepackage[polish]{babel}
\usepackage{lmodern}
\usepackage[a4paper,margin=2.4cm]{geometry}
\usepackage{amsmath}
\usepackage{hyperref}

\title{HF20.3: rozdzielenie błędów technicznych i jakościowych}
\author{Edward Szczypka}
\date{\today}

\begin{document}
\maketitle

\section{Klasyfikacja}

Niech \(H\) oznacza zbiór błędów technicznych kontraktu, a \(Q\) zbiór
niepowodzeń bram jakościowych. Definiujemy:
\[
  \mathrm{serving\_contract\_passed} \iff H=\varnothing,
\]
\[
  \mathrm{publication\_ready} \iff H=\varnothing \land Q=\varnothing.
\]

Jeżeli \(H=\varnothing\) oraz \(Q\neq\varnothing\), wynik jest częściowy:
\[
  \mathrm{partial\_success}=\mathrm{true}.
\]

\section{Opad}

Niepowodzenie bramy opadu nie unieważnia poprawnej siatki czasu ani modeli
PM10, PM2.5 i temperatury. Parametry opadowe otrzymują status
\texttt{experimental}, a pozostałe mogą przejść do lokalnego etapu
serwowania.

\end{document}
